% Options for packages loaded elsewhere
\PassOptionsToPackage{unicode}{hyperref}
\PassOptionsToPackage{hyphens}{url}
\PassOptionsToPackage{dvipsnames,svgnames,x11names}{xcolor}
\documentclass[
  11pt,
]{article}
\usepackage{xcolor}
\usepackage[margin=25mm]{geometry}
\usepackage{amsmath,amssymb}
\setcounter{secnumdepth}{-\maxdimen} % remove section numbering
\usepackage{iftex}
\ifPDFTeX
  \usepackage[T1]{fontenc}
  \usepackage[utf8]{inputenc}
  \usepackage{textcomp} % provide euro and other symbols
\else % if luatex or xetex
  \usepackage{unicode-math} % this also loads fontspec
  \defaultfontfeatures{Scale=MatchLowercase}
  \defaultfontfeatures[\rmfamily]{Ligatures=TeX,Scale=1}
\fi
\usepackage{lmodern}
\ifPDFTeX\else
  % xetex/luatex font selection
\fi
% Use upquote if available, for straight quotes in verbatim environments
\IfFileExists{upquote.sty}{\usepackage{upquote}}{}
\IfFileExists{microtype.sty}{% use microtype if available
  \usepackage[]{microtype}
  \UseMicrotypeSet[protrusion]{basicmath} % disable protrusion for tt fonts
}{}
\makeatletter
\@ifundefined{KOMAClassName}{% if non-KOMA class
  \IfFileExists{parskip.sty}{%
    \usepackage{parskip}
  }{% else
    \setlength{\parindent}{0pt}
    \setlength{\parskip}{6pt plus 2pt minus 1pt}}
}{% if KOMA class
  \KOMAoptions{parskip=half}}
\makeatother
\usepackage{listings}
\newcommand{\passthrough}[1]{#1}
\lstset{defaultdialect=[5.3]Lua}
\lstset{defaultdialect=[x86masm]Assembler}
\usepackage{longtable,booktabs,array}
\usepackage{calc} % for calculating minipage widths
% Correct order of tables after \paragraph or \subparagraph
\usepackage{etoolbox}
\makeatletter
\patchcmd\longtable{\par}{\if@noskipsec\mbox{}\fi\par}{}{}
\makeatother
% Allow footnotes in longtable head/foot
\IfFileExists{footnotehyper.sty}{\usepackage{footnotehyper}}{\usepackage{footnote}}
\makesavenoteenv{longtable}
\setlength{\emergencystretch}{3em} % prevent overfull lines
\providecommand{\tightlist}{%
  \setlength{\itemsep}{0pt}\setlength{\parskip}{0pt}}
\emergencystretch=3em
\AtBeginDocument{\Urlmuskip=0mu plus 2mu}
\lstset{breaklines=true,basicstyle=\small\ttfamily,columns=fullflexible,keepspaces=true}
\setlength{\tabcolsep}{4pt}
\usepackage{iftex}
\ifXeTeX
  \usepackage{newunicodechar}
  \makeatletter
  \newcommand{\DeclareUnicodeCharacter}[2]{%
    \begingroup
    \uccode`\~="#1\relax
    \uppercase{\def\temp{~}}%
    \expandafter\endgroup
    \expandafter\newunicodechar\expandafter{\temp}{#2}%
  }
  \makeatother
\fi
\ifdefined\DeclareUnicodeCharacter
\DeclareUnicodeCharacter{2192}{\ensuremath{\to}}
\DeclareUnicodeCharacter{2190}{\ensuremath{\leftarrow}}
\DeclareUnicodeCharacter{2194}{\ensuremath{\leftrightarrow}}
\DeclareUnicodeCharacter{21D2}{\ensuremath{\Rightarrow}}
\DeclareUnicodeCharacter{21A6}{\ensuremath{\mapsto}}
\DeclareUnicodeCharacter{2212}{\ensuremath{-}}
\DeclareUnicodeCharacter{2264}{\ensuremath{\leq}}
\DeclareUnicodeCharacter{2265}{\ensuremath{\geq}}
\DeclareUnicodeCharacter{2260}{\ensuremath{\neq}}
\DeclareUnicodeCharacter{2113}{\ensuremath{\ell}}
\DeclareUnicodeCharacter{03A3}{\ensuremath{\Sigma}}
\DeclareUnicodeCharacter{03A6}{\ensuremath{\Phi}}
\DeclareUnicodeCharacter{0393}{\ensuremath{\Gamma}}
\DeclareUnicodeCharacter{221E}{\ensuremath{\infty}}
\DeclareUnicodeCharacter{2208}{\ensuremath{\in}}
\DeclareUnicodeCharacter{2209}{\ensuremath{\notin}}
\DeclareUnicodeCharacter{211A}{\ensuremath{\mathbb{Q}}}
\DeclareUnicodeCharacter{211D}{\ensuremath{\mathbb{R}}}
\DeclareUnicodeCharacter{2124}{\ensuremath{\mathbb{Z}}}
\DeclareUnicodeCharacter{2115}{\ensuremath{\mathbb{N}}}
\DeclareUnicodeCharacter{2200}{\ensuremath{\forall}}
\DeclareUnicodeCharacter{2203}{\ensuremath{\exists}}
\DeclareUnicodeCharacter{2286}{\ensuremath{\subseteq}}
\DeclareUnicodeCharacter{00D7}{\ensuremath{\times}}
\DeclareUnicodeCharacter{03B4}{\ensuremath{\delta}}
\DeclareUnicodeCharacter{2202}{\ensuremath{\partial}}
\DeclareUnicodeCharacter{03B2}{\ensuremath{\beta}}
\DeclareUnicodeCharacter{03B1}{\ensuremath{\alpha}}
\DeclareUnicodeCharacter{03C6}{\ensuremath{\phi}}
\DeclareUnicodeCharacter{03B3}{\ensuremath{\gamma}}
\DeclareUnicodeCharacter{00B9}{\ensuremath{^1}}
\DeclareUnicodeCharacter{00B2}{\ensuremath{^2}}
\DeclareUnicodeCharacter{2070}{\ensuremath{^0}}
\DeclareUnicodeCharacter{2080}{\ensuremath{_0}}
\DeclareUnicodeCharacter{2081}{\ensuremath{_1}}
\DeclareUnicodeCharacter{2082}{\ensuremath{_2}}
\fi
\usepackage{bookmark}
\IfFileExists{xurl.sty}{\usepackage{xurl}}{} % add URL line breaks if available
\urlstyle{same}
\hypersetup{
  pdftitle={PFP V: A Practitioner\textquotesingle s Guide to Checked Proposals},
  pdfauthor={Jeremy H. Carroll},
  colorlinks=true,
  linkcolor={Maroon},
  filecolor={Maroon},
  citecolor={Blue},
  urlcolor={Blue},
  pdfcreator={LaTeX via pandoc}}

\title{PFP V: A Practitioner\textquotesingle s Guide to Checked
Proposals}
\author{Jeremy H. Carroll}
\date{2026-10-02 · Version 0.1.0 · Practitioner guide · Author-review
revision}

\begin{document}
\makeatletter
\begin{center}
{\LARGE \@title\par}
\vspace{0.8em}
{\normalsize \@author\par}
\vspace{0.5em}
{\small \@date\par}
\end{center}
\makeatother

\textbf{Code baseline:}
\nolinkurl{52cb5457f6c485ba548cfd292f6012a941ed22cf}

\subsection{Purpose and reader
outcome}\label{purpose-and-reader-outcome}

This guide takes one small inconsistent graph through a complete
session: declare the instance, propose a change, interpret a refusal,
admit a repair, check the result, and restore an earlier state. It is a
practitioner companion to the Proposal Fidelity Protocol (PFP) research
series and its legacy Symbol Game implementations. It adds no new
theorem or measured model-performance result.

A \textbf{proposal} is a candidate supplied by a person, script, or
model. A \textbf{host} is the program that applies the declared checks.
A \textbf{receipt} records what a particular execution checked and
returned. These roles matter because a plausible proposal can be wrong,
while a successful check can establish less than its surrounding prose
suggests.

The first session uses a scripted proposer. No model account, API key,
or live conversation is needed. You should finish able to explain both
an accepted move and the limits of its acceptance.

\subsection{1. Declare the instance before solving
it}\label{declare-the-instance-before-solving-it}

Suppose three events have proposed offsets on one linear time axis:

\begin{itemize}
\tightlist
\item
  Event b is one unit after event a.
\item
  Event c is one unit after event b.
\item
  Event a is one unit after event c.
\end{itemize}

This is a mathematical offset example. There is no wraparound clock or
recurring-day convention. Declare those conventions differently and you
have a different problem.

Represent the events by vertices a, b, c and the statements by directed
edges. The ordered claim vector is

\[
c=(c_{ab},c_{bc},c_{ca})=(1,1,1).
\]

A \textbf{vertex potential} assigns a number \(f_a,f_b,f_c\) to the
vertices. Its edge differences are

\[
df=(f_b-f_a,\ f_c-f_b,\ f_a-f_c).
\]

The differences always sum to zero: each vertex value appears once with
each sign. The supplied claims sum to three. Consequently, no potential
satisfies all three unchanged claims.

For the common series example, declare edge weights \(w=(1,2,3)\). A
weight is the penalty for leaving one unit of discrepancy on that edge.
It is a modeling choice, not evidence that one event is more truthful.
The complete instance therefore includes the vertices, directed edges,
edge order, exact claims, weights, and permitted changes.

\subsection{2. Separate a residual from a changed
claim}\label{separate-a-residual-from-a-changed-claim}

For a chosen potential, the \textbf{residual} is \(r=c-df\). Its
weighted discrepancy is

\[
D_c(f)=|r_{ab}|+2|r_{bc}|+3|r_{ca}|.
\]

The \textbf{quotient minimum} is the smallest such discrepancy over all
potentials:

\[
\Phi_w(c)=\min_f D_c(f).
\]

The notation distinguishes the score at one chosen potential from the
best possible score for unchanged claims. A poor potential can make
\(D_c(f)\) larger without changing \(\Phi_w(c)\).

Take \(f=(0,-2,-1)\). Then \(df=(-2,1,1)\), so \(r=(3,0,0)\) and
\(D_c(f)=3\). For every potential, the residual still sums to three.
Because all weights are at least one,

\[
D_c(f)\geq |r_{ab}|+|r_{bc}|+|r_{ca}|
\geq |r_{ab}+r_{bc}+r_{ca}|=3.
\]

The displayed potential attains this lower bound. This short argument
establishes the minimum for this instance without requiring trust in a
solver.

Two operations now have different meanings. Choosing that potential
concentrates the original discrepancy on one edge; the original cycle
sum remains three. Replacing the claims themselves by \(c'=(-2,1,1)\)
produces a new consistent instance: \(c'=df\) and \(\Phi_w(c')=0\). The
governor demonstration performs this second operation through
\passthrough{\lstinline!step\_cut(proposed\_claims=...)!}. Acceptance
says the changed claims satisfy the declared graph check. It does not
establish that the changed event statement is factually correct.

\subsection{3. Choose the executable
path}\label{choose-the-executable-path}

The supported tutorial in this package is
\href{https://github.com/dbccpage/the_symbol_game/blob/52cb5457f6c485ba548cfd292f6012a941ed22cf/omega_engine_atlas/coda_atlas/omega_engine_symbol_game/005_pfp_practitioners_field_manual/code/quickstart.py}{code/quickstart.py}.
Its Python wrapper sends exact graph inputs to the pinned Rust
calculator. The repository also contains a separate historical root
\passthrough{\lstinline!quickstart.py!}; the commands here intentionally
name the package file.

Use Python 3.10 or newer. The package is self-contained: it carries
copies of the two modules it needs from the repository, the native graph
adapter
\href{https://github.com/dbccpage/the_symbol_game/blob/52cb5457f6c485ba548cfd292f6012a941ed22cf/omega_engine_atlas/coda_atlas/omega_engine_symbol_game/005_pfp_practitioners_field_manual/code/native_functional.py}{\nolinkurl{code/native_functional.py}}
and the Paper 004 graph governor
\href{https://github.com/dbccpage/the_symbol_game/blob/52cb5457f6c485ba548cfd292f6012a941ed22cf/omega_engine_atlas/coda_atlas/omega_engine_symbol_game/005_pfp_practitioners_field_manual/code/l1_symbol_game_governor.py}{\nolinkurl{code/l1_symbol_game_governor.py}}.
Each copy begins with a header that names its source path and commit,
gives the SHA-256 of the source bytes, and lists the local changes.
Apart from its header, the adapter differs from the source baseline only
in import order. The governor copy is later than the baseline, and this
package patches it in two places: it uses the native calculator only
when one is configured and then lets every native refusal propagate,
computing the minimum by exact Python enumeration otherwise (Section 4),
and it refuses a CUT that omits a declared edge (Section 5). Inside a
checkout of the source repository,
\nolinkurl{code/tests/test_vendored_sources.py} compares both copies
with the recorded Git objects; elsewhere that comparison skips and says
why.

The native calculator itself is not part of this package. Build the
\passthrough{\lstinline!solve\_graph\_phi1!} executable from the
reviewed source with the repository's Rust workspace, on any platform:

\begin{lstlisting}
cargo build --locked -j 1 --manifest-path omega_engine/crates_rs/Cargo.toml -p omega_functional --bin solve_graph_phi1
\end{lstlisting}

The following Windows PowerShell setup, run from the repository root,
assumes the reviewed executable already exists at the shown path:

\begin{lstlisting}
$env:OMEGA_FUNCTIONAL_GRAPH_BIN = (Resolve-Path omega_engine/crates_rs/target/debug/solve_graph_phi1.exe).Path
$env:OMEGA_FUNCTIONAL_GRAPH_SHA256 = (Get-FileHash -Algorithm SHA256 $env:OMEGA_FUNCTIONAL_GRAPH_BIN).Hash.ToLowerInvariant()
\end{lstlisting}

Use the actual output path if Cargo has a custom target directory. Other
platforms use their native executable name. The digest pins the
executable bytes used by the adapter; it is not an independent
attestation that those bytes were reproducibly built from the intended
source.

Now, from this package's directory, run (on Windows,
\passthrough{\lstinline!py -3.12!} in place of
\passthrough{\lstinline!python!}):

\begin{lstlisting}
python code/quickstart.py --non-interactive
python code/play_demo.py
\end{lstlisting}

The repository's
\href{https://github.com/dbccpage/the_symbol_game/blob/52cb5457f6c485ba548cfd292f6012a941ed22cf/omega_engine_education/FUNCTIONAL_BOUNDARY.md}{functional-boundary
notes} describe the native setup at the baseline. The transport
implementation is
\href{https://github.com/dbccpage/the_symbol_game/blob/52cb5457f6c485ba548cfd292f6012a941ed22cf/omega_engine_atlas/coda_atlas/omega_engine_symbol_game/005_pfp_practitioners_field_manual/code/native_functional.py}{native\_functional.py}.
A missing executable, missing or mismatched digest, native refusal, or
malformed native response prevents an accepted native result. Treat that
outcome as an environment or evidence issue to resolve before relying on
the host calculation.

\subsubsection{Offline arithmetic when the native calculator is
unavailable}\label{offline-arithmetic-when-the-native-calculator-is-unavailable}

You can still understand and check the finite argument. Run this
ordinary Python snippet, or evaluate it by hand:

\begin{lstlisting}[language=Python]
from fractions import Fraction as Q

claims = (Q(1), Q(1), Q(1))
potential = (Q(0), Q(-2), Q(-1))
edge_differences = (
    potential[1] - potential[0],
    potential[2] - potential[1],
    potential[0] - potential[2],
)
residual = tuple(c - d for c, d in zip(claims, edge_differences))
assert edge_differences == (Q(-2), Q(1), Q(1))
assert residual == (Q(3), Q(0), Q(0))
assert sum(edge_differences) == 0
assert sum(residual) == 3
\end{lstlisting}

This is an arithmetic worksheet. It neither calls the native adapter nor
emits a canonical functional receipt. Its lower-bound explanation is the
argument in Section 2. The package test
\nolinkurl{code/tests/test_offline_worksheet.py} runs this block as
printed. Record a native run as unavailable until the required
calculator can execute.

\subsection{4. Follow the complete
session}\label{follow-the-complete-session}

\href{https://github.com/dbccpage/the_symbol_game/blob/52cb5457f6c485ba548cfd292f6012a941ed22cf/omega_engine_atlas/coda_atlas/omega_engine_symbol_game/005_pfp_practitioners_field_manual/code/play_demo.py}{play\_demo.py}
imports the package's copy of the
\href{https://github.com/dbccpage/the_symbol_game/blob/52cb5457f6c485ba548cfd292f6012a941ed22cf/omega_engine_atlas/coda_atlas/omega_engine_symbol_game/005_pfp_practitioners_field_manual/code/l1_symbol_game_governor.py}{graph
governor}. Its graph uses weights (1,2,2), a legacy variant of the
common series example's (1,2,3). Both have the same unique cheapest
edge, minimum three, and displayed repair. The all-zero potential costs
five in the script and six in the common example; those are scores of
one potential, not the quotient minima.

When \nolinkurl{OMEGA_FUNCTIONAL_GRAPH_BIN} is set, the governor asks
that native calculator, and any refusal or error it returns stops the
step; no Python value is substituted. When it is not set, the governor
computes the quotient minimum by exact Python enumeration of
spanning-tree potentials. Without a configured calculator the session
therefore still runs, but its printed values are Python enumeration
results, not native results. Edge weights must be exact integers or
fractions on both paths; floating-point and Boolean weights are refused.
The package test \nolinkurl{code/tests/test_play_demo.py} checks the
table below on the enumeration path, and again through the native
calculator when pins are configured.

\begin{longtable}[]{@{}
  >{\raggedright\arraybackslash}p{(\linewidth - 4\tabcolsep) * \real{0.3333}}
  >{\raggedright\arraybackslash}p{(\linewidth - 4\tabcolsep) * \real{0.3333}}
  >{\raggedright\arraybackslash}p{(\linewidth - 4\tabcolsep) * \real{0.3333}}@{}}
\toprule\noalign{}
\begin{minipage}[b]{\linewidth}\raggedright
Step
\end{minipage} & \begin{minipage}[b]{\linewidth}\raggedright
Proposed action
\end{minipage} & \begin{minipage}[b]{\linewidth}\raggedright
Expected observation
\end{minipage} \\
\midrule\noalign{}
\endhead
\bottomrule\noalign{}
\endlastfoot
Declare & Create vertices a, b, c and weighted edges & The host has a
finite declared carrier \\
ASSERT & Install claims \(1,1,1\) & The original quotient minimum is
3 \\
CUT & Propose replacement \(5,0,0\) &
\passthrough{\lstinline!DefectIncrease!}, not admitted \\
CUT & Propose replacement \(-2,1,1\) &
\passthrough{\lstinline!ExactRepair!}, admitted, 3 to 0 \\
VERIFY & Request \passthrough{\lstinline!ZERO\_OBSTRUCTION!} & The host
checks its stored claims and admits the check \\
RETRACT & Restore the saved post-ASSERT snapshot & The original claims
and minimum 3 return \\
\end{longtable}

The bad proposal sums to five. Its optimum is five by the same
lower-bound argument, with all discrepancy placeable on the weight-one
edge. The host refuses an increase from three to five. A refused
proposal does not replace the active claims.

The good proposal changes the cheapest edge from one to minus two. It
supplies a consistent replacement. The demonstration saves a snapshot
before either CUT and uses that snapshot's actual identifier when
retracting. Do not guess a snapshot identifier from a printed step
count.

ASSERT installs the starting problem; it is not a claim that every
initialization step decreases discrepancy. The governor's CUT checks
permit non-increase. A tied score can therefore return
\passthrough{\lstinline!IdempotentStationary!}. Strict progress and
finite termination of an arbitrary search require additional conditions;
repeating an admitted stationary move is not a solution strategy.

\subsection{5. Read refusals as
information}\label{read-refusals-as-information}

A refusal identifies a failed check. First determine which check ran and
which input caused it. Then change the relevant proposal or declare a
new instance when the intended problem requires different rules.

\begin{longtable}[]{@{}
  >{\raggedright\arraybackslash}p{(\linewidth - 4\tabcolsep) * \real{0.3333}}
  >{\raggedright\arraybackslash}p{(\linewidth - 4\tabcolsep) * \real{0.3333}}
  >{\raggedright\arraybackslash}p{(\linewidth - 4\tabcolsep) * \real{0.3333}}@{}}
\toprule\noalign{}
\begin{minipage}[b]{\linewidth}\raggedright
Observed result
\end{minipage} & \begin{minipage}[b]{\linewidth}\raggedright
Meaning in this implementation
\end{minipage} & \begin{minipage}[b]{\linewidth}\raggedright
Next operator action
\end{minipage} \\
\midrule\noalign{}
\endhead
\bottomrule\noalign{}
\endlastfoot
\passthrough{\lstinline!CarrierViolation!} & A proposed endpoint is
outside the declared vertices & Correct the endpoint or explicitly
formulate another carrier \\
\passthrough{\lstinline!MoveViolation!} & A proposed edge is not one of
the declared edges & Use a declared relation or revise the instance
specification \\
\passthrough{\lstinline!IncompleteProposal!} & A CUT's replacement
claims omit a declared edge or give one a value that is not an exact
integer or fraction & Supply exactly one exact value for every declared
edge; an omitted edge is never read as a zero claim \\
\passthrough{\lstinline!InvariantViolation!} & A prohibited input, such
as a self-edge, reached the gate & Inspect the exact failed input
condition \\
\passthrough{\lstinline!DefectIncrease!} & The proposed score exceeds
the current score under the selected CUT mode & Compute why it increased
before proposing another change \\
\passthrough{\lstinline!IdempotentStationary!} & A permitted move left
the checked score unchanged & Identify a new argument or record the lack
of progress \\
\passthrough{\lstinline!VerificationFailed!} & The requested obligation
did not pass & Inspect that obligation, rather than treating a different
successful check as a substitute \\
\end{longtable}

The native transport can also fail before a governor verdict exists.
Keep the transport exception and native response, if any, with the run
record. A missing calculation is not a zero discrepancy.

Use integers and \passthrough{\lstinline!Fraction!} inputs for
mathematical values. The tutorial rejects floats and Boolean values at
its mathematical input boundary. In particular, Python's treatment of
\passthrough{\lstinline!True!} as an integer-like value must not turn it
into a declared numerical witness.

\subsection{6. Reproduce a result and inspect its
custody}\label{reproduce-a-result-and-inspect-its-custody}

A useful reproduction record names the source revision, exact input,
command, Python version, native executable digest, exit status, and
output. If files change during a run, identify which version each
observation belongs to. Preserve refused cases alongside admitted ones:
they show whether the gate discriminates.

Run the package checks from this package's directory, one test file per
command:

\begin{lstlisting}
python -m pytest code/tests/test_offline_worksheet.py -q
python -m pytest code/tests/test_play_demo.py -q
python -m pytest code/tests/test_quickstart_native.py -q
\end{lstlisting}

The
\href{https://github.com/dbccpage/the_symbol_game/blob/52cb5457f6c485ba548cfd292f6012a941ed22cf/omega_engine_atlas/coda_atlas/omega_engine_symbol_game/005_pfp_practitioners_field_manual/code/tests/conftest.py}{fixture}
consumes only explicitly supplied \nolinkurl{OMEGA_FUNCTIONAL_GRAPH_BIN}
and \nolinkurl{OMEGA_FUNCTIONAL_GRAPH_SHA256}. It discovers no checkout,
builds no program, and creates no pin. Missing configuration skips, with
that reason, every test that requests native execution. A malformed
digest, unreadable configured executable, or digest mismatch fails those
tests with an explicit message. The worksheet, the offline session, the
refusal, exact-input, parser, and scaffold checks, and the
vendored-source header checks still run without native configuration.
The commands are replay instructions, not a reported fresh pass.

The existing
\href{https://github.com/dbccpage/the_symbol_game/blob/52cb5457f6c485ba548cfd292f6012a941ed22cf/omega_engine_atlas/coda_atlas/omega_engine_symbol_game/005_pfp_practitioners_field_manual/code/paper_033_execution_receipt.json}{execution
receipt},
\href{https://github.com/dbccpage/the_symbol_game/blob/52cb5457f6c485ba548cfd292f6012a941ed22cf/omega_engine_atlas/coda_atlas/omega_engine_symbol_game/005_pfp_practitioners_field_manual/code/cross_domain_symbol_game_receipt.json}{cross-domain
receipt}, and series completion receipt are historical artifacts until
their source and executable bindings are reconciled with a new run. A
SHA-256 digest can bind bytes. It does not establish their truth,
authenticate an unsigned producer, or supply a missing experiment. A
restricted sorted-JSON encoding also needs its own stated contract; it
should not be called full
\href{https://www.rfc-editor.org/rfc/rfc8785}{RFC 8785} conformance
without checking that specification.

\subsection{7. Map a new problem
carefully}\label{map-a-new-problem-carefully}

Before adapting the example, write down what a vertex means, what an
edge value measures, why the edge-difference relation is appropriate,
and which changes are allowed. Supply units and direction conventions.
Decide whether changing a claim is an acceptable correction or would
replace an observation you must preserve.

For example, measured offsets might need an uncertainty model rather
than unrestricted replacement. A dependency relation might be Boolean
rather than an additive quantity. Encoding either as this graph is a
modeling step that requires justification. Renaming vertices to
``currency,'' ``permission,'' or ``quantum state'' does not supply that
justification.

The package's
\href{https://github.com/dbccpage/the_symbol_game/blob/52cb5457f6c485ba548cfd292f6012a941ed22cf/omega_engine_atlas/coda_atlas/omega_engine_symbol_game/005_pfp_practitioners_field_manual/code/multi_domain_symbol_game_bouts.py}{four
labeled graph scenarios} illustrate repeated use of the finite
interface. Their labels do not establish market equilibrium,
authorization security, or a quantum no-go theorem. The governed scripts
make no LLM calls, so hypothetical unconstrained-chat costs in older
records are not a measured comparison.

For model integration, keep the proposer behind the same boundary: parse
the proposed fields, check the declared carrier and operations,
recompute the relevant property, and save the actual response. Consult
\href{https://github.com/dbccpage/the_symbol_game/blob/52cb5457f6c485ba548cfd292f6012a941ed22cf/omega_engine_atlas/coda_atlas/omega_engine_symbol_game/005_pfp_practitioners_field_manual/code/llm_streaming_harness.py}{the
streaming harness} as an implementation example; reviewing its adapter
does not establish a live provider run. The harness accepts integers and
exact rational strings as claim values and refuses a JSON float or
Boolean before any scoring, since
\passthrough{\lstinline!Fraction(0.1)!} is the binary expansion of the
float rather than one tenth
(\nolinkurl{code/tests/test_harness_exact_inputs.py}).

\subsection{8. Connections and next
reading}\label{connections-and-next-reading}

The familiar idea of an untrusted producer supplying evidence to a
checker appears in Necula's
\href{https://doi.org/10.1145/263699.263712}{proof-carrying code}. That
work concerns proofs of a declared code-safety policy. It supplies
context for separating proposal production and evidence checking; this
tutorial does not claim its graph receipt is such a proof.

The minimization and lower-bound reasoning connect to linear-programming
duality; see Boyd and Vandenberghe,
\href{https://doi.org/10.1017/CBO9780511804441}{Convex Optimization},
Chapter 5. The series' \href{../001/paper.pdf}{graph mathematics paper}
develops the finite theory; \href{../004/paper.pdf}{Paper 004} explains
native verification and evidence custody. The
\href{../007/paper.pdf}{curriculum-design guide} turns examples into
lessons, and the \href{../009/paper.pdf}{review protocol} explains how
to audit a claimed result.

This guide supports a finite, inspectable proposal--check workflow. Its
practical success criterion is that the reader can state the original
instance, explain the failed proposal, distinguish residual
concentration from claim replacement, and identify exactly what a fresh
receipt records. General semantic correctness and model quality need
their own evidence.

\subsection{9. Limitations}\label{limitations}

The guide works one three-edge instance by hand and describes one
scripted session. The package tests check the session table on the
governor's exact Python enumeration path. For this revision the
native-requesting tests were also run against a
\passthrough{\lstinline!solve\_graph\_phi1!} executable built on Windows
(x86\_64-pc-windows-msvc, Rust 1.98.0) with
\passthrough{\lstinline!cargo build --locked -j 1!} from the
\passthrough{\lstinline!omega\_functional!} crate, whose source is
unchanged since the code baseline; with that executable and its SHA-256
(\nolinkurl{2b0379557993246cacb1ed8964f9f096964019ecc1c2fb781e6b18f2ca6c8a65})
pinned, all 75 package tests passed and none was skipped. The dated
review receipt records that run. The executable is not part of the
deposit, so another reader must build and pin their own. When a native
executable is configured the governor uses it and propagates its
refusals; otherwise it uses the enumeration, so a printed value does not
by itself show which calculator produced it. The legacy scenario and
streaming scripts are implementation examples, not measured comparisons.

This paper assigns no authorization and changes no jurisdiction. It
carries no operational authority or institutional standing.

\subsection{Verification for this
revision}\label{verification-for-this-revision}

The package verifier (\passthrough{\lstinline!verify\_deposit.py!})
checks the manifest and the two-way file census, and
\passthrough{\lstinline!tools/rebuild\_release.py --check!} checks the
release archive against the manifests. The package tests and lint are
executed by the dated review receipt recorded with this revision; that
receipt lists each command, tool version, exit status, and tested scope.
Historical receipts in this package keep their original dates and scope
and are not current execution evidence.

\subsection{References}\label{references}

\begin{itemize}
\tightlist
\item
  Necula, G. C. (1997).
  \href{https://doi.org/10.1145/263699.263712}{Proof-carrying code}. In
  \emph{Proceedings of the 24th ACM SIGPLAN-SIGACT Symposium on
  Principles of Programming Languages (POPL '97)}, pp.~106--119. ACM.
  DOI: 10.1145/263699.263712.
\item
  Boyd, S., and Vandenberghe, L. (2004).
  \href{https://doi.org/10.1017/CBO9780511804441}{Convex Optimization}.
  Cambridge University Press, Chapter 5. DOI:
  \nolinkurl{10.1017/CBO9780511804441.}
\item
  Rundgren, A., Jordan, B., and Erdtman, S. (2020).
  \href{https://www.rfc-editor.org/rfc/rfc8785}{JSON Canonicalization
  Scheme (JCS)}. RFC 8785 (Informational), RFC Editor, June 2020. DOI:
  10.17487/RFC8785.
\end{itemize}

\end{document}
